\documentclass[11pt]{article}
\usepackage[a4paper, margin=2cm]{geometry}
\usepackage{graphicx}
\usepackage{amsmath}
\usepackage{enumitem}

\begin{document}

\begin{enumerate}

\item Fig. 1.1 shows a skydiver falling vertically through the air.

\begin{center}
\includegraphics[width=0.6\textwidth]{fig1_1.png} \\
\textit{Fig. 1.1}
\end{center}

In the first part of the fall, her speed increases and her acceleration decreases. \\
In the second part of the fall, her speed is constant.

\begin{enumerate}[label=(\alph*)]
    \item On Fig. 1.2 sketch the speed-time graph for the skydiver.
    On your graph, mark one point where the speed of the skydiver is increasing with an A and one point where the speed of the skydiver is constant with a B.
    
    \begin{center}
    \includegraphics[width=0.5\textwidth]{fig1_2.png} \\
    \textit{Fig. 1.2}
    \end{center}

    \item Explain how the graph shows that the acceleration decreases as the speed increases.

    \item During the first part of the fall, there is a resultant vertical force acting downwards on the skydiver.
    \begin{enumerate}[label=(\roman*)]
        \item One of the vertical forces acting on the skydiver is her weight.
        State the name of the other vertical force that acts on the skydiver.
        
        \item Explain why the resultant vertical force eventually becomes zero.
    \end{enumerate}

    \item At one instant, the vertical force on the skydiver is 400 N downwards.
    At the same instant, the wind causes an additional horizontal force of 100 N to the right to act on the skydiver.
    Draw a vector diagram to determine the resultant of the 400 N vertical force and the 100 N horizontal force. Place an arrow on all of the forces to show their directions.
    Determine the magnitude (size) of this resultant force and its direction to the vertical.
\end{enumerate}

\item 
\begin{enumerate}[label=(\alph*)]
    \item State the principle of moments.

    \item An airline passenger wishes to check the weight of a suitcase that he carries as hand luggage onto an aeroplane. \\
    He uses a uniform plank of wood, pivoted at its centre, as shown in Fig. 2.1.
    
    \begin{center}
    \includegraphics[width=0.8\textwidth]{fig2_1.png} \\
    \textit{Fig. 2.1}
    \end{center}
    
    The plank is balanced with the suitcase on one side of the pivot and three bags of sugar, each of mass 2.0 kg, on the other side.
    The distances are shown on Fig. 2.1.
    
    \begin{enumerate}[label=(\roman*)]
        \item Calculate the weight of one bag of sugar.
        
        \item The maximum weight of a suitcase that can be carried onto the aeroplane is 67 N.
        Determine whether the weight of the suitcase exceeds the maximum allowed.
        Show a calculation in your answer.
    \end{enumerate}

    \item A tall bus is tested for stability.
    Fig. 2.2 shows the bus on a slope.
    The centre of gravity of the bus is marked.
    
    \begin{center}
    \includegraphics[width=0.5\textwidth]{fig2_2.png} \\
    \textit{Fig. 2.2}
    \end{center}
    
    \begin{enumerate}[label=(\roman*)]
        \item State what is meant by the 'centre of gravity' of an object.
        
        \item When the slope is made very steep, the bus falls over by rotating about point P.
        Explain why the bus falls over.
    \end{enumerate}
\end{enumerate}

\item A battery, a pulley and a motor are used to lift a load as shown in Fig. 3.1.

\begin{center}
\includegraphics[width=0.4\textwidth]{fig3_1.png} \\
\textit{Fig. 3.1}
\end{center}

\begin{enumerate}[label=(\alph*)]
    \item Describe the transfers between energy stores that occur as the load is lifted.

    \item The efficiency of the motor, pulley and load system is less than 100\%.
    \begin{enumerate}[label=(\roman*)]
        \item By comparing the input energy and the useful output energy, explain why the efficiency is less than 100\%.
        
        \item Explain how the principle of conservation of energy applies in lifting the load.
    \end{enumerate}

    \item The input power to the motor is 15 W. The motor is used for 20 s. The efficiency of the motor is 60\%.
    Calculate the energy supplied to the load.
\end{enumerate}

\item 
\begin{enumerate}[label=(\alph*)]
    \item Evaporation of water from the surface of the skin causes cooling.
    \begin{enumerate}[label=(\roman*)]
        \item Describe, using ideas about particles, how evaporation causes cooling.
        
        \item State one difference between evaporation and boiling.
    \end{enumerate}

    \item When a refrigerator is switched on, cooling coils placed at the top of the space inside the refrigerator become cold. This causes a convection current which cools the air inside the refrigerator. The refrigerator is shown in Fig. 4.1.
    
    \begin{center}
    \includegraphics[width=0.5\textwidth]{fig4_1.png} \\
    \textit{Fig. 4.1}
    \end{center}
    
    \begin{enumerate}[label=(\roman*)]
        \item Explain how the cooling coils cause a convection current in the air inside the refrigerator.
        
        \item The food in the refrigerator is initially at a temperature of $20.0^{\circ}\text{C}$. The food has a mass of 3.6 kg and a specific heat capacity of $3000 \text{ J}/(\text{kg}^{\circ}\text{C})$.
        Calculate the final temperature of the food after 160000 J of thermal energy is removed from it.
    \end{enumerate}
\end{enumerate}

\item 
\begin{enumerate}[label=(\alph*)]
    \item Fig. 5.1 is a diagram showing the arrangement of air particles as a longitudinal wave passes through them.
    
    \begin{center}
    \includegraphics[width=0.8\textwidth]{fig5_1.png} \\
    \textit{Fig. 5.1}
    \end{center}
    
    \begin{enumerate}[label=(\roman*)]
        \item On Fig. 5.1, mark the centre of a compression with the letter C, and mark the centre of a rarefaction with the letter R.
        
        \item Describe the difference between a compression and a rarefaction.
    \end{enumerate}

    \item In a ripple tank, a water wave is produced by a wooden bar moving up and down on the surface of water.
    \begin{enumerate}[label=(\roman*)]
        \item The wooden bar makes 45 complete oscillations in 1.0 minute.
        Calculate the frequency of the wave produced.
        
        \item The frequency of the water wave is increased by moving the wooden bar up and down more quickly.
        State what happens to the speed and what happens to the wavelength of the wave produced.
        
        \item The crests of the water wave move into the shallow region shown in Fig. 5.2.
        
        \begin{center}
        \includegraphics[width=0.5\textwidth]{fig5_2.png} \\
        \textit{Fig. 5.2}
        \end{center}
        
        On Fig. 5.2, draw the crests in the shallow region.
    \end{enumerate}

    \item Describe what is meant by the diffraction of a water wave.
\end{enumerate}

\item A student sets up a circuit to determine the resistance of a length of wire. \\
The circuit contains a battery of unknown e.m.f., a length of a wire used to make a resistor X, an ammeter, a voltmeter and a variable resistor R.

\begin{enumerate}[label=(\alph*)]
    \item Fig. 6.1 shows part of the circuit diagram.
    
    \begin{center}
    \includegraphics[width=0.5\textwidth]{fig6_1.png} \\
    \textit{Fig. 6.1}
    \end{center}
    
    \begin{enumerate}[label=(\roman*)]
        \item On Fig. 6.1, complete the circuit diagram by adding one voltmeter and one ammeter in suitable places to allow the determination of the resistance of X.
        
        \item Explain how X and R act as a variable potential divider.
    \end{enumerate}

    \item The student determines the resistance of the resistor X for five different lengths of the wire making it. The lengths of wire range from 20 cm to 60 cm. The type of wire and the cross-sectional area of the wire are kept constant.
    Fig. 6.2 shows a graph of the results.
    
    \begin{center}
    \includegraphics[width=0.6\textwidth]{fig6_2.png} \\
    \textit{Fig. 6.2}
    \end{center}
    
    \begin{enumerate}[label=(\roman*)]
        \item State the relationship between the resistance of the wire and the length of the wire.
        
        \item Calculate the current in a 90 cm length of the wire when there is a potential difference (p.d.) of 9.0 V across it.
        Show your working.
        
        \item The p.d. across the wire making the resistor X is kept constant for all the measurements of resistance.
        Describe the relationship between the current in the wire and the length of the wire.
    \end{enumerate}
    
    \item State how the resistance of a wire depends upon the cross-sectional area of the wire.
\end{enumerate}

\item Fig. 7.1 shows an alternating current (a.c.) power supply connected to a transformer.

\begin{center}
\includegraphics[width=0.6\textwidth]{fig7_1.png} \\
\textit{Fig. 7.1}
\end{center}

\begin{enumerate}[label=(\alph*)]
    \item Explain how an alternating current in the primary coil produces an alternating output voltage.

    \item A student uses a voltmeter set on a 0-10 V range to measure the input and output voltages.
    She obtains the values shown in Table 7.1.
    
    \begin{center}
    \textbf{Table 7.1} \\
    \begin{tabular}{|c|c|}
    \hline
    \textbf{input voltage / V} & \textbf{output voltage / V} \\
    \hline
    1.2 & 2.4 \\
    2.3 & 4.6 \\
    4.8 & 9.6 \\
    6.4 & no reading \\
    \hline
    \end{tabular}
    \end{center}
    
    \begin{enumerate}[label=(\roman*)]
        \item Suggest why no output voltage reading is obtained with this voltmeter when the input voltage is 6.4 V.
        
        \item The number of turns on the primary coil is 48.
        Calculate the number of turns on the secondary coil.
    \end{enumerate}

    \item The student uses an oscilloscope to display an alternating output voltage from the transformer.
    Fig. 7.2 shows the front of the oscilloscope before it is connected to the transformer.
    
    \begin{center}
    \includegraphics[width=0.6\textwidth]{fig7_2.png} \\
    \textit{Fig. 7.2}
    \end{center}
    
    When the oscilloscope is connected to the output of the transformer, a trace representing the alternating output voltage is displayed on the screen.
    \begin{enumerate}[label=(\roman*)]
        \item On Fig. 7.2, draw a trace representing an alternating output voltage on the screen.
        
        \item Describe how to use the trace to measure the maximum value of the output voltage.
    \end{enumerate}
\end{enumerate}

\item Plutonium-239 (${}^{239}_{94}\text{Pu}$) is an isotope that is used as the fuel in some nuclear reactors.

\begin{enumerate}[label=(\alph*)]
    \item State the names of the types of particles found in a nucleus of plutonium-239, and state how many there are of each type.

    \item Fig. 8.1 shows the nuclear fission process that occurs within the fuel rods of the nuclear reactor.
    
    \begin{center}
    \includegraphics[width=0.6\textwidth]{fig8_1.png} \\
    \textit{Fig. 8.1}
    \end{center}
    
    \begin{enumerate}[label=(\roman*)]
        \item Explain how the fission process produces a chain reaction.
        
        \item Explain how control rods are used to increase and decrease the rate of the chain reaction in a nuclear reactor.
    \end{enumerate}

    \item Plutonium-239 decays by the emission of an alpha particle ($\alpha$-particle).
    State two differences between an alpha particle and a beta particle ($\beta$-particle).
    
    \item Alpha particles from the radioactive source are detected in a cloud chamber or with a spark counter.
    \begin{enumerate}[label=(\roman*)]
        \item Draw a labelled diagram of either a cloud chamber or a spark counter.
        Label the position of the radioactive source with an S.
        
        \item State what causes the tracks in a cloud chamber or state what causes the sparks in a spark counter.
    \end{enumerate}
\end{enumerate}

\item 
\begin{enumerate}[label=(\alph*)]
    \item The life cycle of a star begins with a large cloud of dust and gas which collapses.
    Five later stages of the life cycle of a very massive star are:
    
    \begin{center}
    black hole \hspace{1cm} protostar \hspace{1cm} red supergiant \hspace{1cm} stable star \hspace{1cm} supernova
    \end{center}
    
    Place these stages in Table 9.1 in the order in which they occur.
    
    \begin{center}
    \textbf{Table 9.1} \\
    \begin{tabular}{|p{5cm}|}
    \hline
    \textbf{earlier time} \\
    \hline
    cloud of dust and gas \\
    \hline
    \rule{0pt}{0.75cm} \\
    \hline
    \rule{0pt}{0.75cm} \\
    \hline
    \rule{0pt}{0.75cm} \\
    \hline
    \rule{0pt}{0.75cm} \\
    \hline
    \rule{0pt}{0.75cm} \\
    \hline
    \textbf{later time} \\
    \hline
    \end{tabular}
    \end{center}

    \item The original collapse of the cloud of dust and gas that formed the Sun was caused by an inward force.
    \begin{enumerate}[label=(\roman*)]
        \item State the name of the inward force.
        
        \item Further collapse is prevented by an outward force. The Sun will remain in the stable stage of its life cycle for a few billion years.
        Describe what causes the outward force.
    \end{enumerate}

    \item One of the first supernovas ever observed is known as SN185. It was formed from the explosion of a star in the Milky Way galaxy.
    The remnants of SN185 are at a distance of 8200 light-years from Earth.
    \begin{enumerate}[label=(\roman*)]
        \item State what is meant by a 'light-year'.
        
        \item State the time that passed between the explosion that formed SN185 and the electromagnetic radiation from the explosion reaching Earth.
        
        \item A recently observed supernova is SN2014J.
        The remnants of SN2014J are 12 million light-years from Earth, outside the Milky Way.
        There is no redshift seen in the electromagnetic radiation from the remnants of SN185 but a large redshift is seen in the electromagnetic radiation from the remnants of SN2014J.
        Explain this difference.
    \end{enumerate}

    \item Most of the atoms found in the early Universe were hydrogen and helium.
    The Universe now contains atoms of heavier elements.
    Explain how the heavier elements are formed.
\end{enumerate}

\end{enumerate}
\end{document}
